\justifying
\NSFCBodyText

\subsubsection{年度研究计划}

研究期限为两年，与实验室开放课题管理要求一致。

\textbf{第一年（2027 年 1 月—12 月）。} 完成公开数据合规审查与划分核对；在 SUST 上训练场景文字编码器并在 RUST 上测试；实现无文字通道的图文对比基线与 OCR—翻译流水线；搭建 Recall@$K$、MRR 与消融记录脚本；提交年度进展报告。阶段目标是确认 $f_s$ 可用，并得到 B0--B2 的可复现分数。

\textbf{第二年（2028 年 1 月—12 月）。} 完成门控融合与视觉枢纽训练；在 Multi30k-Distant 上报告单语与跨语言检索；完成 200--400 题理解协议与专家核对；撰写并投稿 SCI 论文；按实验室要求提交结题报告。阶段目标是完成 H1--H3 的检验，并形成署名符合实验室第一单位要求的成果。

\subsubsection{预期研究结果}

方法上，形成维语场景文字增强的多模态融合与跨语言检索评价协议。实验上，报告完整模型、三类基线与四项消融。成果上，力争发表中科院三区及以上 SCI 论文 1 篇；论文由申请人所在单位与教育部重点实验室共同署名，并以实验室为第一单位。不预先承诺专利、平台上线或学生数据实验。

\subsubsection{学术交流与合作计划}

拟在课题执行期内参加 1 次国内多模态或民族语言信息处理学术会议，汇报检索与消融结果。与实验室的合作以合同、年度报告和成果署名为准，不安排未落实的国际合作经费。差旅与会务按实验室实报实销科目列支。预算总额 5 万元，按申请人确定列为劳务费 3.5 万元、业务费 1.5 万元。业务费对应测试费 0.8 万元、会议费/差旅费 0.4 万元、出版物/文献/信息费 0.3 万元。劳务费能否按实验室科目报销，以立项合同和实报实销审定为准。立项后首期使用 50\% 经费，中期报告达到优或良后再使用其余 50\%。
